\section*{Overview}

Virtual tree is the standard way to avoid touching an entire tree when a query only mentions a small set of marked
vertices. It builds the smallest tree that still contains all marked vertices and all paths between them.

\section*{When to Use It}

Use it when:

\begin{itemize}
  \item the base graph is a tree,
  \item each query marks a relatively small subset of vertices,
  \item the needed computation depends only on paths between marked nodes.
\end{itemize}

\section*{Core Idea}

Take the marked vertices, sort them by Euler-tour entry time, add the LCAs of consecutive marked vertices, sort and
deduplicate again, then connect them with a stack using ancestor relationships.

The resulting compressed tree has only \(O(k)\) vertices for a query with \(k\) marked nodes.

\section*{Key Insight}

On a tree, the only extra vertices you need are LCAs. Nothing else is necessary to preserve the path structure between
the marked nodes. That is why the virtual tree stays small.

\section*{Worked Problem}

\paragraph{Problem.}
Each query gives \(k\) special cities in a country road tree. Compute some DP over the union of paths connecting those
cities, for example the minimum number of checkpoints to separate all special cities from each other.

\paragraph{Why virtual tree fits.}
The answer only depends on the Steiner tree of the special cities. The virtual tree is exactly that structure in a form
that is easy to DP over.

\section*{Correctness Intuition}

Sorting by Euler time means ancestor intervals are nested correctly. The stack connects each node to the nearest earlier
node that is its ancestor in the compressed set, which is exactly the parent it has in the minimal path-preserving tree.

\section*{Complexity Analysis}

If a query marks \(k\) nodes, building the virtual tree costs \(O(k \log k)\) with ordinary sorting, plus \(O(1)\) or
\(O(\log n)\) per LCA depending on the chosen LCA structure.

\section*{Implementation}

The sample code assumes you already have:

\begin{itemize}
  \item Euler entry times \texttt{tin},
  \item an \texttt{lca(u, v)} routine,
  \item an \texttt{is\_ancestor(u, v)} test.
\end{itemize}

It returns the compressed node list and the virtual-tree adjacency.

\section*{Common Pitfalls}

\begin{itemize}
  \item Forgetting to insert LCAs before building the stack.
  \item Building edges in the original-tree depth order instead of Euler order.
  \item Reusing adjacency from one query without clearing it.
\end{itemize}

\section*{Variants / Extensions}

\begin{itemize}
  \item Weighted virtual edges storing original-tree distance.
  \item Query-by-query DP on marked nodes.
  \item Combining with centroid or HLD preprocessing around the base tree.
\end{itemize}

\section*{Practice Problems}

\begin{itemize}
  \item DP on marked nodes of a tree.
  \item Count or optimize over the union of marked-node paths.
  \item Queries where \(k \ll n\) and the full tree is too expensive per query.
\end{itemize}

\section*{References}

\begin{itemize}
  \item \href{https://oiwiki.com/graph/virtual-tree/}{OI Wiki: Virtual Tree}.
  \item \href{https://codeforces.com/catalog}{Codeforces Catalog}.
  \item \href{https://usaco.guide/plat/VT?lang=cpp}{USACO Guide: Virtual Tree}.
\end{itemize}
